\section{Competency Questions and Query-Based Evaluation}
\label{sec:ontoksd-competency-questions}

OntoKSD is examined through competency questions concerning competence
specifications, instructional use, assessment provenance, evidence,
proficiency interpretation, and learner-related records. The questions
make these information needs testable through identified SPARQL
operations. Their execution examines whether the represented information
can be retrieved under specified conditions; it does not, by itself,
establish the educational validity of a competence or an assessment
judgment.

The revised collection distinguishes \emph{conceptual question families}
from \emph{executable query views}. It retains 42 historical families,
identified as CQ1--CQ42, implemented by 58 root-level query files:
57 read-only \texttt{SELECT} queries and one optional update, CQ39b.
These counts exclude fixture-loading scripts and verification queries
in the test subdirectories. Some views are complementary, while others
are equivalent retrievals retained for traceability. The number of
files is therefore not a count of independent capabilities or satisfied
requirements.

\subsection{Identifier Stability after the Qualification Review}

The original family identifiers are retained following the
post-qualification revision. Renumbering the executable files as a
single consecutive sequence would obscure their correspondence with
the earlier question set, the thesis tables, and the repository
artifacts. Alphabetic suffixes identify operational views within a
family: CQ21a and CQ21b, for example, distinguish evidence-origin
provenance from participation provenance. The nested suffixes CQ20c1
and CQ20c2 retain the existing identifiers of focused realization and
context views under CQ20. A suffix does not introduce an additional
functional requirement.

The revision also supplies explicit implementations for CQ14--CQ17,
which were listed in the question set but absent as individual files
from the inspected directory. Their addition closes that implementation
omission; it is not evidence that those files had already been executed
in GraphDB. The repository identifier map records every file, family,
operation, and equivalent view so that historical references remain
resolvable without imposing a misleading consecutive file count.

\subsection{Competency Questions and Requirement Mapping}

Table~\ref{tab:traceability-cq-fr} states the information need of each
family and its principal requirement mapping. FR14 applies across the
collection. The mapping identifies intended examination, rather than
asserting complete requirement satisfaction from the existence of a
query. The operational files and their result semantics are documented
in the repository catalog.

\begingroup
\small
\begin{longtable}{p{0.08\linewidth} p{0.70\linewidth} p{0.12\linewidth}}
\caption{Historical competency-question families and principal functional requirement mappings.}
\label{tab:traceability-cq-fr} \\
\toprule
\textbf{CQ} & \textbf{Information need in the revised implementation} & \textbf{FRs} \\
\midrule
\endfirsthead
\toprule
\textbf{CQ} & \textbf{Information need in the revised implementation} & \textbf{FRs} \\
\midrule
\endhead
\bottomrule
\endfoot
\multicolumn{3}{l}{\textbf{A. Competence Specifications and Structure}} \\
\midrule
CQ1 & Which direct constituents specify a selected competence, including knowledge, skills, dispositions, attitudes, and knowledge--level pairings? & FR1, FR3 \\
CQ2 & Which \textit{KnowledgeSkill} pairings and recorded justification verbs are associated with a selected competence? & FR2 \\
CQ3 & Which knowledge element and required skill level constitute each populated \textit{KnowledgeSkill} pairing? & FR2 \\
CQ4 & Which specification-pattern value is associated with a competence? & FR1 \\
CQ5 & Which atomic or composite structure is declared by a structure value or represented through class membership? & FR1, FR4 \\
CQ6 & Which direct component relationships and integration descriptions are recorded? & FR4 \\
CQ7 & Which direct specialization relationships and narrower-scope descriptions are recorded? & FR4 \\
CQ8 & Which records are represented as reusable competence specifications, whole-competence judgments, or constituent judgments? & FR1, FR11 \\
\midrule
\multicolumn{3}{l}{\textbf{B. Instructional Associations and Conditions}} \\
\midrule
CQ9 & Which instructional units are associated with a competence through activation? & FR5 \\
CQ10 & Which functions, roles, modes, constraints, or generic qualifiers characterize an activation? & FR5 \\
CQ11 & Which competences are explicitly positioned as prerequisites of instructional units? & FR5 \\
CQ12 & Which competences are explicitly positioned as instructional targets? & FR5 \\
CQ13 & Which modes qualify an activation? & FR5 \\
CQ14 & Which instructional contribution roles qualify an activation? & FR5 \\
CQ15 & Which constraints, classification types, and descriptions qualify an activation? & FR5 \\
CQ16 & Which function, roles, modes, and constraints are jointly associated with the same activation? & FR5 \\
CQ17 & Which independent instructional statements, collaboration settings, and planned acquisition contexts describe instructional units, including units without activation records? & FR5, FR6 \\
\midrule
\multicolumn{3}{l}{\textbf{C. Evidence and Assessment Records}} \\
\midrule
CQ18 & Which evidence supports whole-competence or constituent judgments? & FR10, FR11 \\
CQ19 & Which evidence types classify the collected material? & FR10 \\
CQ20 & Which occurrences ground evidence, which assessment specifications do they realize, and in which acquisition contexts do they occur? & FR10 \\
CQ21 & How do whole-competence judgments connect to evidence-origin occurrences and to the occurrences of their assessment participations? & FR10, FR11 \\
CQ22 & How does a whole-competence judgment connect its competence, profile, supporting evidence, participation, and occurrence? & FR11 \\
CQ23 & How does a constituent judgment connect its specification target, evidence, participation, profile, and optional parent judgment? & FR11 \\
CQ24 & Which learners, assessors, optional scores, derived judgments, and occurrence-origin evidence are associated with assessment participations? & FR10, FR11 \\
CQ25 & Which judgment dates and evidence-generation or validation dates are available across assessment occasions? & FR11, FR13 \\
\midrule
\multicolumn{3}{l}{\textbf{D. Proficiency Interpretation}} \\
\midrule
CQ26 & Which proficiency level and interpretation scheme are assigned to a whole-competence judgment? & FR12 \\
CQ27 & Which proficiency level and interpretation scheme are assigned to a constituent judgment? & FR12 \\
CQ28 & In which participation occurrence is each proficiency judgment grounded? & FR11, FR12 \\
CQ29 & Which proficiency schemes are explicitly associated with assessment occurrences? & FR12 \\
CQ30 & Which numeric scales are used by proficiency schemes, with which available bounds and increment? & FR12 \\
CQ31 & Which intervals, endpoint policies, scales, and corresponding levels form a numeric interpretation scheme? & FR12 \\
CQ32 & Which overall proficiency is already assigned to a whole-competence judgment? The implemented view CQ32a repeats CQ26; the originally envisaged synthesis procedure is not implemented. & Supplementary \\
\midrule
\multicolumn{3}{l}{\textbf{E. Curricular Relationships and Learner-Oriented Retrieval}} \\
\midrule
CQ33 & Which course specifications and instructional activations associate a competence with curricular or instructional units? & FR5 \\
CQ34 & Which external alignments identify internal and external endpoints, correspondence types, identifiers, and available justifications? & FR8 \\
CQ35 & In which distinct curricular or instructional contexts is the same competence reused? & FR7 \\
CQ36 & How do activation qualifications differ across uses of that competence? & FR5, FR7 \\
CQ37 & Which whole-competence judgments, levels, and schemes are associated with a profile? & FR12, FR13 \\
CQ38 & Which constituent judgments can be traced to a profile through participation or a parent judgment? & FR11, FR13 \\
CQ39 & Which eligible judgments fall below a selected curricular minimum, which candidates can optionally be materialized, and which learning needs are explicitly recorded? & Supplementary \\
CQ40 & Which mandatory expectations lack judgments or have comparable below-minimum judgments, and which direct constituents lack records? Constituent threshold propagation remains disabled by default. & Supplementary \\
CQ41 & Which whole-competence and constituent judgments, specification targets, levels, and schemes derive from a participation? & FR10--FR13 \\
CQ42 & Which competence-targeted learning needs retain documented comparison bases, which comparisons repeat CQ39a, and which curricular inclusions have no judgment in a profile? & Supplementary \\
\end{longtable}
\endgroup

The revised wording follows the implemented operations. In particular,
CQ7 concerns direct specialization rather than unspecified semantic
relations; CQ17 retrieves instructional statements independently of
activation; and CQ32a retrieves an existing attribution rather than
calculating a synthesis. These distinctions prevent broader ambitions
of the initial question set from being reported as implemented results.

Several views deliberately overlap. CQ13--CQ15 inspect individual
qualifiers, whereas CQ16 retains their common activation identifier in
a joint view. Multiple modes or constraints can produce combinations
of values, so rows must not be counted as distinct activations. CQ26
and CQ32a are equivalent whole-competence proficiency retrievals.
CQ39a, CQ40b, and CQ42b apply the same below-minimum comparison and
must yield equivalent results under the same parameters, reference
time, and dataset. Such equivalence is a regression condition, not
three independent demonstrations of requirement coverage.

\subsection{Datasets, Parameters, and Reproducible Execution}

The evaluation distinguishes three sources of instance data:
competence and instructional specifications derived from the CSP and
curricular examples; an explicitly fictional assessment and proficiency
example; and dedicated synthetic fixtures for query behavior absent
from those sources. Planning records do not, by themselves, provide
learner evidence, assessment participations, or proficiency judgments.
Synthetic records must therefore remain identifiable and excluded
from empirical counts of CSP outcomes.

Each execution must record the query version, repository, selected
dataset or graphs, ontology and vocabulary versions, inference
configuration, parameters, and exported result. Ordinary query files
do not impose a dataset through \texttt{FROM}; endpoint configuration
determines which graphs are visible. The dedicated test views specify
their synthetic graph explicitly. Graph separation identifies the test
source but does not ensure exclusion from a default dataset configured
to include named graphs.

The profile/course queries now provide executable defaults referring
to \texttt{Profile\_Learner\_01} and \texttt{Course\_FormalLanguages}
in the fictional proficiency example. Both IRIs must be replaced for
a different case. This explicit selection avoids assuming enrollment
or applying every course expectation to every profile. CQ1 and CQ2
retain \texttt{C02} as their selected competence example.

The three proficiency scripts supply, in order, the interpretation
vocabulary, assessment-related instances and example score
classification, and a curricular expectation with a local comparison
policy. Their updates are separate from read-only CQ execution.
CQ39b is likewise an update and is not included in a batch of
retrieval operations. It materializes comparison candidates only when
intentionally executed; previously stored candidates are not
invalidated automatically by later evidence.

For temporal comparison, the distributed files use \texttt{NOW()}.
A reproducible experiment fixes a typed reference instant and retains
the resulting query. The local regression run uses
\texttt{2026-11-01T12:00:00Z} as that analytical instant, not as a claim
about when an assessment occurred or a GraphDB execution was performed.

\subsection{Empty Results and Synthetic Test Cases}

An empty answer is evaluated against a stated expectation. It can
indicate an unmatched example parameter, an absent relation, an
unavailable dataset, a deliberately excluded record, or a policy that
has not been enabled. It does not alone demonstrate a query defect,
zero proficiency, or lack of learner competence. Optional labels,
contexts, and descriptors must be distinguished from mandatory joins.
Where a field is optional, its absence should remain visible rather
than remove an otherwise matching record.

CQ11 has no prerequisite activations in the examined CSP cycles. A
dedicated synthetic fixture therefore tests three explicit prerequisite
associations and control activations representing targets or having
no declared function. It includes the same competence as a target in
one unit and a prerequisite in another, and an association expressed
through the inverse unit link. This tests contextual classification
and retrieval without asserting that the proposed prerequisites are
pedagogically necessary. Course-to-course prerequisites remain outside
this question.

CQ27 is empty in the original proficiency example because that example
does not assign levels to its constituent judgments. A separate fixture
provides two constituent-specific interpretations under textual rubrics:
conceptual knowledge is classified as consolidated, while construction
skill is classified as developing. Both use the same response material,
but their criteria and justifications differ; no overall score is copied
to the constituents. A deliberately incomplete third record has a level
but no recorded scheme and must remain visible with that column unbound.
Controls without a level and of the whole-competence judgment class are
excluded. These controls test retrieval and missing information; they
are not additional validated educational conclusions.

CQ20c requires an occurrence type, while realization and context are
optional. CQ30 requires only the scheme--scale association; its optional
bounds, increment, and labels cannot establish the existence of that
association. Its revised label filters accept Portuguese and untagged
literals within the optional blocks. English-only entities consequently
remain retrievable with blank label fields. This is neither ranked
language fallback nor a guarantee of one row per scale. The earlier
empty result is not attributed to missing data or a query defect solely
from its subsequent successful execution: the recorded diagnostic
established that the associations were present in the default dataset,
but did not establish the cause of the earlier execution outcome.

In the fictional proficiency example, the competence has a recorded
judgment at ordinal position~3 and a selected curricular expectation
at position~4. The comparison views can therefore return a below-minimum
candidate when their remaining eligibility conditions hold. By contrast,
CQ40a and CQ42c are expected to be empty for that competence because a
judgment exists. A low recorded result and an absent judgment are
different cases. CQ40c can still identify an unrecorded direct
\textit{KnowledgeSkill} judgment without treating its nested knowledge
component as automatically lacking evidence.

\subsection{Comparison Policy and Supplementary Scope}

The comparison views require a mandatory curricular specification, an
explicitly identified common proficiency scheme, dated judgments, and
unambiguous ordinal category positions. The course specification uses
\texttt{dcterms:references} to identify its expected scheme, while the
judgment uses \textit{interpretedUsingScheme}. This is a documentary
convention of the analytical example, not an additional OWL inference.
Both compared levels must belong to that scheme. Equal ordinals across
different schemes do not establish semantic equivalence.

Within the selected scheme, the policy considers the latest dated,
non-expired record at the reference instant. Future or expired records,
conflicting same-date levels, and multiple incompatible assignments
cannot justify choosing a convenient lower value. A newer unclassified
record blocks use of an older classified result within the implemented
selection conditions. Ordinal comparison establishes an ordering of
categories, not a measured distance between them.

The curricular minimum describes expected depth of coverage. Its
comparison with individual performance is an explicitly adopted
analytical policy, rather than a learner passing threshold or an
ontology-level diagnosis. CQ39a identifies candidates without requiring
stored \textit{LearningNeed} instances. CQ39c and CQ42a instead retrieve
recorded needs and their documented basis, which need not describe the
latest current judgment. CQ40a concerns mandatory inclusions without a
judgment; CQ42c also includes non-mandatory curricular inclusions and
reports their status.

CQ40d remains disabled. The ontology does not assign a whole-competence
minimum to every constituent, and the available examples do not justify
that propagation. Enabling its flag would introduce an external policy
assumption, not activate an established constituent requirement. The
synthetic CQ27 levels do not remove this limitation. Constituent-specific
gap analysis would require justified expectations at that granularity
or a separately defended comparison policy. Likewise, CQ32a provides
no synthesis or automatic aggregation procedure; that original
supplementary ambition remains unimplemented.

\subsection{Verification Evidence and Coverage Limits}

The local revision checks parsed all 58 root-level files and executed
the 57 SELECT queries on the local ontology and fictional proficiency
example. Additional artificial data exercised the new CQ14--CQ17 views
and the independent-statement case. The dedicated CQ11 and CQ27
fixtures returned their three expected associations or records. The
local checks also confirmed equivalent views, expected empty results,
and idempotent materialization on a disposable graph. They used
RDFLib/SPARQL~1.1 without fetching ontology imports or performing OWL
reasoning. The repository validation report records input and query
hashes and the fixed analytical instant.

These are local regression results, not a claim that every revised file
has been executed successfully in GraphDB. Interactive GraphDB checks
reported during revision are not a complete execution log for the final
collection. In particular, newly supplied views and revised query hashes
require corresponding runs and exports before a final GraphDB success
claim can be made. Logical consistency with imported modules requires
separate reasoner checks. The assessment criteria defined in
Section~\ref{subsec:ontoksd-requirement-verification} remain applicable.

The question set provides evidence toward the functional baseline with
the following limits:
\begin{itemize}
 \item CQ17 now retrieves independent instructional statements even
 without activation, addressing a retrieval aspect of FR6. It does not
 establish that the recorded statement captures every action or
 deliverable of the source task.
 \item CQ3 retrieves the required skill level but does not explicitly
 project its framework-specific reference scheme. FR2 therefore also
 requires examination of level typing and the imported vocabulary.
 \item CQ6 and CQ7 already retrieve integration and specialization
 descriptions; CQ34 already retrieves correspondence type and
 justification. The remaining checks concern presence, scope, and
 adequacy of those descriptions, rather than an assumed absence of
 these fields from the query projections.
 \item FR9 and the criteria-related part of FR11 require examination
 of expected products, observations, and applicable assessment criteria.
 Recovering evidence or a rubric reference does not establish criterion
 adequacy or that the evidence supports the recorded conclusion.
 \item CQ30 and CQ31 concern numeric structures. CQ27's synthetic
 rubric example demonstrates retrieval of qualitative interpretations,
 but not exhaustive validation of qualitative schemes or their use in
 the empirical CSP data.
 \item Provenance and profile paths require separate integrity checks
 for agreement about learner, assessment occasion, constituent
 membership, and evidence origin. Retrieval alone does not enforce
 those path constraints.
\end{itemize}

FR14 is examined through explicit query artifacts and recorded execution
conditions. It is not satisfied merely by counting files or non-empty
answers. Requirement-level conclusions must distinguish support,
partial support, lack of support, and absence of evaluation according
to the relevant checks in
Table~\ref{tab:ontoksd-requirement-checks}.

The FR-to-CQ view in Appendix~\ref{ap:OntoKSD-FR-CQ} must retain these
family identifiers, the revised CQ17 mapping, and the same coverage
limits. Terminology is consolidated in
Appendix~\ref{ap:glossary-ontoksd}. Together, the family table, executable
catalog, parameterized tests, and execution records provide the basis
for examining the implemented architecture without conflating
representational capability, query behavior, and educational validity.
